%% ma: L11-B27-0015
%% lop: 11
%% chuong: 7
%% bai: 27
%% dang: 11.27.11
%% loai: TLN
%% muc: VD
%% dap_an: "5"
%% nguon: "(NBV)-ĐỀ SỐ 13-MỖI NGÀY 1 ĐỀ THI 2025.pdf (D13, MỖI NGÀY 1 ĐỀ - Đề số 13), Phần III Câu 2"
%% kien_thuc: "Bài 27 (hình chóp cụt đều), Bài 22 (góc giữa hai đường thẳng), Bài 24 (định lí cosin lớp 10)"
%% trang_thai: cho-duyet
%% ly_do: "Dạng mới 11.27.11 (côsin góc giữa đường chéo và cạnh đáy của khối chóp cụt đều) chờ giáo viên duyệt"
%% ghi_chu: "Đáp án gốc 5 đúng (cos = 2/3, a + b = 5, kiểm số: cạnh bên bằng căn 3, chiều cao khối bằng 1, đường chéo bằng 3). Đã bỏ ảnh khay đá minh họa; không vẽ hình vì đề không đưa hình toán học. Lưu ý đề gốc viết chiều cao mặt bên là căn 2 cm (đường cao của hình thang cân mặt bên)."
\begin{ex}
(THPT Lý Thường Kiệt -- Hà Nội 2025) Một khay đá gồm $6$ ngăn nhỏ có dạng là các hình chóp cụt với miệng đáy là hình vuông. Ta đo được độ dài cạnh đáy nhỏ, cạnh đáy lớn lần lượt bằng $1$ cm và $3$ cm và chiều cao mặt bên bằng $\sqrt2$ cm. Tính $a+b$ biết côsin góc giữa đường chéo của viên đá với cạnh đáy của viên đá có dạng $\dfrac ab$ ($a$, $b$ là hai số nguyên tố cùng nhau).
\shortans{5}
\loigiai{Gọi hình chóp cụt đều có đáy lớn $ABCD$ (cạnh $3$) và đáy nhỏ $A'B'C'D'$ (cạnh $1$). Mặt bên là hình thang cân có hai đáy $3$ và $1$, đường cao $\sqrt2$, nên cạnh bên $AA'=\sqrt{2+1}=\sqrt3$.

Hình chiếu $H$ của $A'$ lên $(ABCD)$ nằm trên $AC$ với $AH=\dfrac{AC-A'C'}2=\dfrac{3\sqrt2-\sqrt2}2=\sqrt2$, nên $A'H=\sqrt{3-2}=1$ và $A'C=\sqrt{HC^2+A'H^2}=\sqrt{(2\sqrt2)^2+1}=3$.

Góc giữa đường chéo $A'C$ và cạnh đáy $CB$ là $\widehat{A'CB}$. Mặt bên $ABB'A'$ là hình thang cân nên $A'B^2=AA'^2+AB\cdot A'B'=3+3=6$. Trong tam giác $A'CB$ có $A'C=CB=3$, $A'B=\sqrt6$, nên
\[\cos\widehat{A'CB}=\dfrac{9+9-6}{2\cdot3\cdot3}=\dfrac23.\]
Vậy $a=2$, $b=3$ và $a+b=5$.}
\end{ex}
